% Zero_Decoherence_Check.tex
% Critical cross-reference audit: Which argument does the paper actually use
% for zero decoherence, and is it consistent?
% Date: 2026-03-23

\documentclass[11pt,a4paper]{article}
\usepackage[utf8]{inputenc}
\usepackage{amsmath,amssymb,amsthm}
\usepackage{geometry}
\usepackage{xcolor}
\usepackage{tcolorbox}
\usepackage{booktabs}
\geometry{margin=2.5cm}

\title{Zero-Decoherence Prediction: Cross-Reference Audit\\
\large Which argument does the paper actually use?}
\author{DFD Research Programme --- Cross-Reference Check}
\date{2026-03-23}

\begin{document}
\maketitle

%=============================================================================
\section{The Three Possible Arguments}
%=============================================================================

There are three logically distinct arguments that could support a
``no intrinsic decoherence from $\psi$'' prediction. They have different
premises, different scopes, and different vulnerabilities.

\begin{description}
\item[Argument (1): Branch-by-branch sourcing.]
$\psi$ tracks the actual matter configuration on each quantum branch
separately. On each branch, matter is in a definite position, and
$\psi$ responds to that definite configuration via uniqueness of
the elliptic BVP. There is one $\psi$ per branch; no conflict between
branches; no decoherence. This requires a Leray--Schauder-type or
monotone-operator uniqueness argument to guarantee that each branch
uniquely determines its $\psi$.

\item[Argument (2): $\langle\rho\rangle$ sourcing with linear response.]
$\psi$ responds to the quantum expectation value $\langle\hat\rho\rangle$
but the response is ``linear enough'' that coherence is preserved.
For a \emph{nonlinear} field equation, this argument is \textbf{wrong}:
$\psi[\langle\rho\rangle] \neq \langle\psi[\rho]\rangle$ when the
map $\rho \mapsto \psi$ is nonlinear. The nonlinear deviation would
actually \emph{cause} decoherence (the superposition sees a different
$\psi$ than either branch).

\item[Argument (3): $S_\psi/\hbar \gg 1$ suppresses quantum fluctuations
of $\psi$ itself.]
The action of the $\psi$ field is vastly super-Planckian, so $\psi$
does not have independent quantum fluctuations. This is a separate
claim about the \emph{quantum nature of $\psi$ itself}, not about
its response to quantum matter. It is true but does not address the
Page--Geilker problem.
\end{description}

%=============================================================================
\section{What the Paper Actually Says}
%=============================================================================

\subsection{Section 14 (section\_openproblems.tex), lines 30--50}

The paper presents the Penrose paradox resolution as follows
(Sec.~14.1, ``Quantum Superpositions and the Penrose Paradox''):

\begin{tcolorbox}[colback=red!5!white, colframe=red!60!black,
  title=Paper's Actual Text (lines 31--48)]
\textbf{Line 31:} ``The resolution follows from the linearity of the
source equation''

\textbf{Lines 36--41:} For $|\Psi\rangle = c_A|A\rangle + c_B|B\rangle$:
\begin{enumerate}
\item The source density is
$\rho = |c_A|^2\rho_A + |c_B|^2\rho_B$ (quantum expectation value)
\item The $\psi$ field responds to this weighted average
\item No ``branched geometry'' exists; there is one $\psi$ field for
the system
\end{enumerate}

\textbf{Lines 44--48:} DFD predicts no intrinsic decoherence from $\psi$-field
at current experimental scales; standard unitary QM evolution unless
environmental decoherence dominates.
\end{tcolorbox}

\subsection{Section 14.2 (same file), line 65}

The paper also states (in the UV completion subsection):
\begin{quote}
``Classical $\psi$ by design: The action scales as
$S_\psi \sim (M_{\mathrm{Planck}}/a_\star)^2 \gg \hbar$,
ensuring quantum fluctuations of $\psi$ are negligible. The field
doesn't need quantization.''
\end{quote}

\subsection{Section 3 (section\_wellposedness.tex), line 559--563}

The well-posedness section states:
\begin{quote}
``The semi-classical regime (quantum matter on classical $\psi$ background)
is well-defined. Full quantization of $\psi$ is unnecessary:
$S_\psi \sim (M_P/a_\star)^2 \gg \hbar$, ensuring quantum fluctuations
are negligible.''
\end{quote}

%=============================================================================
\section{Diagnosis}
%=============================================================================

\begin{tcolorbox}[colback=yellow!10!white, colframe=orange!80!black,
  title=\textbf{FINDING: The paper uses Argument (2), which is problematic}]

The paper explicitly states:
\begin{itemize}
\item $\rho = |c_A|^2\rho_A + |c_B|^2\rho_B$ --- this is the
\textbf{quantum expectation value} $\langle\hat\rho\rangle$.
\item ``The $\psi$ field responds to this weighted average'' ---
this is the M{\o}ller--Rosenfeld semiclassical prescription.
\item The paper then \emph{also} invokes Argument (3)
($S_\psi \gg \hbar$) in Sec.~14.2 and Sec.~3.6, but this addresses
a \emph{different} question (whether $\psi$ itself has quantum
fluctuations, not how it responds to quantum matter).
\end{itemize}

\textbf{The paper does NOT use Argument (1).} There is no
branch-by-branch picture, no reference to Leray--Schauder, and no
invocation of the uniqueness theorem to establish that each quantum
branch determines its own $\psi$.
\end{tcolorbox}

%=============================================================================
\section{Why This Matters}
%=============================================================================

\subsection{The Page--Geilker problem}

The $\langle\hat\rho\rangle$ sourcing prescription was experimentally
tested by Page and Geilker (1981). A radioactive decay triggered one
of two macroscopic mass configurations near a Cavendish balance. If
$\psi$ (or, in their case, the gravitational field) responded to
$\langle\hat\rho\rangle$, the balance would have shown the
\emph{weighted average} of both configurations. Instead, it showed
one definite configuration --- the one corresponding to the actual
decay outcome.

The paper's current formulation (lines 36--41) explicitly uses the
$\langle\hat\rho\rangle$ prescription and is therefore directly
contradicted by Page--Geilker.

\subsection{The nonlinearity problem}

Even setting aside Page--Geilker, the paper's claim of ``linearity
of the source equation'' (line 31) is misleading. The DFD field
equation
\[
\nabla\cdot[\mu(|\nabla\psi|/a_\star)\nabla\psi] = -\frac{8\pi G}{c^2}\rho
\]
is \textbf{linear in $\rho$} (the right-hand side) but
\textbf{nonlinear in $\psi$} (the left-hand side). Therefore the
map $\rho \mapsto \psi[\rho]$ is nonlinear. This means:
\[
\psi[|c_A|^2\rho_A + |c_B|^2\rho_B]
\;\neq\;
|c_A|^2\,\psi[\rho_A] + |c_B|^2\,\psi[\rho_B].
\]
The $\psi$ field that responds to $\langle\hat\rho\rangle$ is
\emph{not} the average of the branch-by-branch $\psi$ fields.
This nonlinear mismatch would actually \emph{cause} decoherence in
the $\langle\hat\rho\rangle$ picture, contradicting the paper's
``no decoherence'' prediction.

\subsection{The $S_\psi \gg \hbar$ argument is orthogonal}

Argument (3) correctly establishes that $\psi$ does not develop
independent quantum fluctuations. But this is about whether $\psi$
\emph{itself} is quantum, not about how it \emph{responds} to quantum
matter. Even a perfectly classical field can cause decoherence if it
mediates distinguishable interactions between branches. The
$S_\psi \gg \hbar$ statement is true and relevant for the UV completion
discussion, but it does not resolve the Penrose/Page--Geilker problem.

%=============================================================================
\section{Leray--Schauder Check}
%=============================================================================

\textbf{Question:} Does section\_wellposedness.tex reference the
Leray--Schauder theorem by name?

\textbf{Answer: No.} A thorough search of section\_wellposedness.tex
(and all other .tex files in the paper) finds \textbf{zero} occurrences
of ``Leray'' or ``Schauder''. The uniqueness results in Section 3 are
based on:
\begin{itemize}
\item \textbf{Strict monotonicity} of the flux operator
$\mathbf{a}(\xi) = \mu(|\xi|)\xi$ (Theorem 3.2)
\item \textbf{Strict convexity} of the energy functional (proof via
calculus of variations)
\item \textbf{Crandall--Liggett} semigroup theory (for the parabolic
extension, Theorem 3.8)
\end{itemize}

The uniqueness statement (Theorem 3.2) is: ``If the flux operator
$\mathbf{a}(\xi)$ is strictly monotone [strict inequality in (A4)],
then the weak solution of Theorem 3.1 is unique.''

This monotone-operator uniqueness \emph{could} support Argument (1)
--- if the paper had connected it to the quantum branching picture.
Each branch has a definite $\rho$; the uniqueness theorem guarantees
a unique $\psi$ for each branch; therefore $\psi$ tracks branches, not
averages. But the paper does not make this connection.

%=============================================================================
\section{Consistency Verdict}
%=============================================================================

\begin{tcolorbox}[colback=red!5!white, colframe=red!70!black,
  title=\textbf{INCONSISTENCY IDENTIFIED}]

\textbf{Status:} The paper's zero-decoherence prediction (Sec.~14.1)
\textbf{implicitly assumes $\langle\hat\rho\rangle$ sourcing}
(Argument 2), which is:
\begin{enumerate}
\item Experimentally falsified by Page--Geilker (1981),
\item Internally inconsistent with the nonlinear field equation
(would actually predict decoherence, not zero decoherence),
\item Not supported by the uniqueness theorems in Section 3
(which could support Argument 1, but the paper does not invoke them
for this purpose).
\end{enumerate}

\textbf{Strength of the claim as written:} WEAK. The prediction of
``no intrinsic decoherence'' is stated but the supporting argument
(response to $\langle\hat\rho\rangle$) undermines rather than
supports it.

\textbf{However:} The \emph{correct} argument (Argument 1:
branch-by-branch sourcing via uniqueness) is available within the
paper's own mathematical framework. The uniqueness theorem
(Theorem 3.2, strict monotonicity) guarantees that each branch's
$\rho$ uniquely determines a $\psi$. This would:
\begin{itemize}
\item Resolve Page--Geilker (each branch sees its own $\psi$),
\item Make the ``no intrinsic decoherence'' prediction rigorous,
\item Be consistent with the nonlinear field equation (no averaging
needed).
\end{itemize}
\end{tcolorbox}

%=============================================================================
\section{Recommended Fix}
%=============================================================================

Replace Sec.~14.1 lines 31--41 with a branch-by-branch argument:

\begin{tcolorbox}[colback=green!5!white, colframe=green!60!black,
  title=Proposed Replacement Text]
\textbf{Why DFD resolves this paradox.}
In DFD, there is \textit{one} flat $\mathbb{R}^3$ with \textit{one}
scalar field $\psi$. The resolution follows from the
\textbf{uniqueness theorem} (Theorem~3.2): for any matter
configuration $\rho$, the strictly monotone flux operator guarantees
a unique solution $\psi[\rho]$.

For a quantum superposition
$|\Psi\rangle = c_A|A\rangle + c_B|B\rangle$:
\begin{enumerate}
\item On branch $A$, matter is at definite positions with density
$\rho_A$; by Theorem~3.2, this uniquely determines $\psi_A$.
\item On branch $B$, matter has density $\rho_B$; uniqueness gives
$\psi_B$.
\item There is no ``branched geometry'' --- each branch has flat
$\mathbb{R}^3$ with a unique scalar field. The quantum superposition
is of matter states, not of spacetimes.
\item The $\psi$-field does not cause decoherence because it does not
distinguish branches \emph{beyond} what the matter already
distinguishes. (The field is a functional of the matter
configuration, not an independent degree of freedom that could
carry which-path information.)
\end{enumerate}
This is the branch-by-branch picture, consistent with the
Page--Geilker experiment (1981) and with standard quantum mechanics.
The classical action $S_\psi \sim (M_P/a_\star)^2 \gg \hbar$
(Sec.~14.2) further ensures that $\psi$ has no independent quantum
fluctuations.
\end{tcolorbox}

%=============================================================================
\section{Cross-Reference Summary}
%=============================================================================

\begin{table}[h]
\centering
\caption{Cross-reference status of decoherence-related claims.}
\begin{tabular}{p{4cm}p{3cm}p{5cm}}
\toprule
\textbf{Claim} & \textbf{Location} & \textbf{Status} \\
\midrule
$\rho = |c_A|^2\rho_A + |c_B|^2\rho_B$ &
Sec.~14.1, line 38 &
\textcolor{red}{PROBLEM: This is $\langle\hat\rho\rangle$ sourcing} \\
No intrinsic decoherence &
Sec.~14.1, line 46 &
\textcolor{orange}{Correct prediction, wrong justification} \\
$S_\psi \gg \hbar$ &
Sec.~14.2, line 65; Sec.~3.6, line 561 &
\textcolor{green!60!black}{True, but addresses different question} \\
Uniqueness theorem &
Sec.~3, Theorem 3.2 &
\textcolor{green!60!black}{Available but not connected to Sec.~14.1} \\
Leray--Schauder &
Not referenced &
\textcolor{orange}{Not needed; monotone operator theory suffices} \\
Branch-by-branch &
Not stated &
\textcolor{red}{MISSING: Should be the core argument} \\
\bottomrule
\end{tabular}
\end{table}

\vspace{1em}
\noindent\textbf{Bottom line:} The paper has the mathematical tools
(Theorem~3.2) to make the zero-decoherence prediction rigorous via
Argument (1), but currently uses Argument (2) which is both wrong and
experimentally falsified. Fixing this requires rewriting approximately
10 lines in Sec.~14.1 to invoke the branch-by-branch picture grounded
in the uniqueness theorem. The prediction itself (``no intrinsic
decoherence'') is correct --- only the justification needs repair.

\noindent\textbf{See also:}
\texttt{Page\_Geilker\_Resolution.tex} in the same directory, which
provides a full rigorous derivation of the branch-by-branch argument.

\end{document}
